\section{Công đoạn 2: Biểu diễn tri thức bản thể học (ontology) và đồ thị tri thức lâm sàng}

\subsection{Khảo sát các hệ thống bản thể học và đồ thị tri thức y sinh}
Khi xử lý văn bản tự do, hệ thống dễ đánh mất ngữ cảnh giải phẫu và làm đứt gãy quan hệ logic giữa các thực thể. Để khắc phục điều này, tri thức cần được cấu trúc hóa bằng bản thể học (ontology) và đồ thị tri thức (Knowledge Graph -- KG). Mục này khảo sát các hệ thống danh mục chuẩn hóa quốc tế và ba công trình tiêu biểu.

\paragraph{Các hệ thống danh mục chuẩn hóa quốc tế (UMLS, SNOMED-CT, RadLex).}
Đây là các từ điển chuẩn hóa thực thể y sinh lớn nhất hiện nay. SNOMED-CT cung cấp hơn 350.000 khái niệm lâm sàng, trong khi RadLex chuyên biệt cho thuật ngữ chẩn đoán hình ảnh. Tuy nhiên, các hệ thống này chỉ mang tính phân cấp tĩnh (\emph{is-a hierarchy}) và thiếu các luật suy diễn lâm sàng gắn với điều kiện ràng buộc thực tế của ca bệnh, chẳng hạn tuổi bệnh nhân, độ di lệch hay bệnh nền.

\paragraph{Biểu diễn tri thức bản thể học bệnh lý xương (Dao \& Ly, 2025) \cite{dao2025ontology}.}
Công trình đặt nền móng cho việc mô hình hóa các tổn thương cơ xương khớp bằng ngôn ngữ bản thể học OWL và hệ luật Semantic Web Rule Language (SWRL). Ưu điểm chính là khả năng suy luận logic vị từ chính xác cho các ca bệnh cơ học. Hạn chế lớn nhất là hệ thống khép kín theo mô hình hệ chuyên gia cổ điển, không giao tiếp được với mô hình ngôn ngữ lớn để đọc hiểu các bệnh án lâm sàng tự do từ thực tế bệnh viện.

\paragraph{Tự động xây dựng đồ thị tri thức lâm sàng bằng nhiều mô hình ngôn ngữ lớn (Das và cộng sự, 2026) \cite{das2025clinicalkg}.}
Công trình đề xuất quy trình trích xuất các bộ ba quan hệ $\langle h, r, t \rangle$ từ báo cáo lâm sàng bằng 3 LLM phối hợp thẩm định chéo. Điểm mạnh của cách làm này là bao quát dữ liệu nhanh, nhưng đồ thị sinh ra vẫn có thể chứa các cạnh nhiễu và không có khung bản thể học chuẩn để ràng buộc các quy tắc an toàn phẫu thuật.

\paragraph{Học đa nguồn tăng cường tri thức bệnh án (KnowGen, 2024) \cite{knowgen2024}.}
Nghiên cứu công bố trên tạp chí \emph{Information Fusion} cho thấy việc kết hợp đồ thị tri thức y khoa vào quá trình mã hóa hồ sơ bệnh án điện tử (EHR) giúp tăng rõ rệt độ nhạy chẩn đoán so với chỉ xử lý chuỗi văn bản.

\subsection{Bảng tiêu chí so sánh các phương pháp biểu diễn tri thức và đồ thị lâm sàng}

Bảng~\ref{tab:comparison_ontology} đối chuẩn các công trình biểu diễn tri thức theo bốn tiêu chí: mức độ bao phủ bệnh lý xương, cơ chế chuẩn hóa thực thể, khả năng suy luận bằng luật hình thức và khả năng tích hợp với LLM.

\begin{table}[H]
\begin{singlespace}
\centering
\fontsize{7.5pt}{9.2pt}\selectfont
\setlength{\tabcolsep}{2pt}
\renewcommand{\arraystretch}{1.15}
\begin{tabular}{|l|l|c|l|l|l|}
\hline
\textbf{Công trình} & \textbf{Nguồn tri thức} & \textbf{Bệnh xương} & \textbf{Chuẩn hóa thực thể} & \textbf{Hệ luật} & \textbf{Tích hợp với LLM} \\
\hline
UMLS / SNOMED-CT & Danh mục quốc tế & \ding{55} & Mã CUI / SCTID & \ding{55} Phân cấp tĩnh & \ding{55} Tra cứu rời rạc \\
\hline
Dao \& Ly (2025) \cite{dao2025ontology} & Bệnh học xương & \ding{51} & OWL, logic mô tả & \ding{51} Luật SWRL & \ding{55} Hệ chuyên gia đóng \\
\hline
Das et al. (2026) \cite{das2025clinicalkg} & Báo cáo lâm sàng, EHR & \ding{55} & LLM kiểm tra chéo & \ding{55} Đồ thị phẳng & \ding{51} LLM tự sinh đồ thị \\
\hline
KnowGen (2024) \cite{knowgen2024} & EHR, đồ thị y khoa & \ding{55} & Ghép vector đặc trưng & \ding{55} & \ding{51} Hợp nhất trên EHR \\
\hline
\textbf{Đề xuất (6 tầng)} & Phác đồ AAOS/WHO & \ding{51} & Chuẩn hóa 6 tầng & \ding{51} Chống chỉ định, cờ đỏ & \ding{51} Dẫn đường cho RAG \\
\hline
\end{tabular}
\caption{So sánh các phương pháp biểu diễn tri thức và đồ thị lâm sàng}
\label{tab:comparison_ontology}
\end{singlespace}
\end{table}

\subsection{Phân tích khoảng trống nghiên cứu (research gaps) ở Công đoạn 2}
Từ Bảng~\ref{tab:comparison_ontology}, có thể xác định ba hạn chế chính.

\paragraph{Khoảng trống 2.1 -- Đứt gãy giữa bản thể học hình thức và ngôn ngữ tự nhiên.}
Công trình tiên phong của Dao \& Ly (2025) \cite{dao2025ontology} đã xây dựng được hệ luật bản thể học xương chặt chẽ, nhưng đây là một hệ thống đóng: nó không xử lý được văn bản bệnh án tự do và không trả lời được câu hỏi bằng ngôn ngữ tự nhiên của bác sĩ.

\paragraph{Khoảng trống 2.2 -- Thiếu ràng buộc quy tắc an toàn trong đồ thị tri thức tự động.}
Các phương pháp sinh đồ thị tự động bằng LLM như Das và cộng sự (2026) \cite{das2025clinicalkg} tạo ra các đồ thị phẳng, thiếu chiều sâu ngữ nghĩa. Các đồ thị này không có các nút phân tầng (tầng bệnh học, tầng phẫu thuật, tầng chống chỉ định) và không có cơ chế ngăn chặn các cạnh quan hệ giả tạo.

\paragraph{Khoảng trống 2.3 -- Bỏ sót các cờ đỏ bệnh lý ác tính trong chấn thương.}
Các đồ thị lâm sàng hiện nay chỉ biểu diễn chấn thương cơ học, bỏ qua mối liên hệ sinh bệnh học giữa gãy xương bệnh lý (\emph{Pathological Fracture}) và u xương ác tính.
